% paper/sections/experiments.tex

\section{Experimental Setup}
\label{sec:experiments}

\subsection{Training Conditions}

We compare four conditions evaluated on the same held-out test distribution:

\textbf{PID baseline.}
A three-axis proportional-integral-derivative controller with gains
$K_p = 20$, $K_i = 0.1$, $K_d = 10$, anti-windup limit $8$\,rad/s,
and goal threshold $0.8$\,m (wider than the $0.5$\,m RL threshold
to account for integral oscillation near the goal).
The PID requires no training and is evaluated directly on both the
training and held-out test distributions to quantify classical control
fragility.

\textbf{Naive SAC.}
SAC trained on fixed nominal physics with no randomisation ($|\Xi| = 1$).
Tests whether RL alone, without distribution coverage, produces
transfer-robust policies.

\textbf{Uniform DR.}
SAC trained with physics parameters sampled uniformly from the full
training ranges (Table~\ref{tab:dr_ranges}) at every episode reset.
This is the standard domain randomisation approach~\cite{tobin2017domain}.

\textbf{CDR (ours).}
SAC trained with performance-gated range expansion as described in
Section~\ref{sec:cdr}. Begins from the narrow CDR start ranges in
Table~\ref{tab:dr_ranges} and expands toward the Uniform DR ranges
based on rolling success rate.

\subsection{Training Protocol}

Each SAC condition is trained for $1 \times 10^6$ environment steps
across three independent random seeds (seeds 0, 1, 2), yielding nine
main experiments.
All conditions share identical SAC hyperparameters (Table~\ref{tab:hyperparams})
and differ only in the physics distribution during training.
Training uses a single NVIDIA T4 GPU on Google Colab and requires
approximately 4.5 hours per run.
Observations are normalised online using running mean and variance
(VecNormalize~\cite{stable_baselines3}), with the statistics saved
alongside each trained model and reloaded at evaluation time.

\subsection{Evaluation Protocol}

All policies are evaluated on a \emph{held-out test distribution}
whose parameter ranges are strictly wider than those seen during training
(Table~\ref{tab:dr_ranges}, rightmost column).
No policy was trained on or fine-tuned with test-distribution data;
all results represent zero-shot transfer.

We run 100 evaluation episodes per seed per condition using deterministic
policy rollout ($\arg\max$ action selection).
Goal distance threshold is $0.5$\,m for all SAC conditions.
We report four metrics:

\begin{itemize}
\item \textbf{Success rate}: fraction of episodes reaching the goal
  within the distance threshold.
\item \textbf{Mean reward}: mean undiscounted episode return $\pm$
  standard deviation across seeds.
\item \textbf{Energy per step}: mean absolute thruster command
  $\frac{1}{T}\sum_t \|a_t\|_1 / 4 \in [0, 1]$,
  proportional to battery energy consumption.
\item \textbf{Mean final distance}: mean goal distance at episode end
  $\pm$ standard deviation.
\end{itemize}

Statistical tests use Welch's $t$-test (unequal variance) with
Cohen's $d$ effect size reported throughout.
Bootstrap 95\% confidence intervals (1000 resamples) are reported
where sample size ($n = 3$ seeds) limits statistical power.

\subsection{Sensor Noise Extension}
\label{sec:master_setup}

We repeat the three-condition SAC comparison with sensor noise added
to the DR parameter space:
position noise $\sigma_p \in [0,\,0.10]$\,m,
velocity noise $\sigma_v \in [0,\,0.05]$\,m/s,
and angular velocity noise $\sigma_\omega \in [0,\,0.02]$\,rad/s.
These represent realistic IMU and DVL sensor degradation.
Results are reported in Section~\ref{sec:noise_results}.

\subsection{Implementation}

All environments use Gymnasium~\cite{gymnasium} with the MuJoCo~\cite{mujoco}
physics backend.
Policies are implemented using Stable-Baselines3~\cite{stable_baselines3}.
Code and trained models are available at:
\url{https://github.com/Itzlimon22/rl_robotics}.